\SourceLeaf{348}

\begin{center}
[III.]

\textbf{ON THE THEORY OF CONIC SECTIONS.}
\end{center}

The investigations deal first with the ellipse, but they can be extended
to the hyperbola with slight modifications.

\Needspace{8\baselineskip}\begin{center}
[1.]
\end{center}

First let $m$, $n$ be two real numbers, $\rho$ a quantity of the
form $\cos\varphi + i\sin\varphi$, so that $\rho$ runs through all its values
as $\varphi$ increases from $0^0$ to $360^0$.

Let
\[
r = m\rho + \dfrac{n}{\rho},
\]
and let $R$ be the point representing the complex quantity $r$; $R$ is thus
a variable point on an ellipse whose principal axes $m+n$, $m-n$ lie on
the lines of the real and imaginary numbers, while its
centre coincides with the origin.

The four values of $\rho$, namely $\alpha$, $\beta$, $\gamma$, $\delta$,
correspond to the four values of $r$, namely $a$, $b$, $c$, $d$, or the four
points $A$, $B$, $C$, $D$.

\Needspace{7\baselineskip}Finally set
\[
\dfrac{-\alpha\beta\gamma m^3 + (\alpha+\beta+\gamma) m m n
  - \delta(\alpha\beta+\alpha\gamma+\beta\gamma) m n n + \delta n^3}
  {(1-\alpha\beta\gamma\delta) m n} = t
\]
and call the point to which the complex number $t$ refers $T$.

It is easy to see that if three of the four points $A$, $B$, $C$, $D$
are fixed and the fourth traverses the conic section, $T$ traverses a straight
line.

\Needspace{9\baselineskip}A simple calculation gives
\[
t - a = -\dfrac{(\alpha\beta m - n)(\alpha\gamma m - n)(m - \alpha\delta n)}
  {\alpha(1-\alpha\beta\gamma\delta) m n}
\]
\[
= -\dfrac{(a-b)(a-c)}{a-d} \cdot
  \dfrac{\alpha\beta\gamma(\alpha-\delta)\left(mm -
  \left(\alpha\delta+\dfrac{1}{\alpha\delta}\right) m n + n n\right)}
  {(\alpha-\beta)(\alpha-\gamma)(1-\alpha\beta\gamma\delta) m n}.
\]
